\section{Validation across arrivals and model deployments}
\label{app:e5e6}
The arrival study holds engine calibration fixed while changing arrivals; the 14B study remeasures
lifecycle timing for a 14B deployment. Both use the same 2100\,s window,
deadlines, four local slots per GPU, and scenario ledger. The original predictions precede
physical execution, and repeats share a prediction; the additional 30\,s
contrast below is retrospective. Table~\ref{tab:e5e6}
reports policy cells, excluding full-capacity anchors.

\subsection{Fixed calibration on additional arrival windows}
The arrival study retains Qwen2.5-7B-Instruct/SGLang 0.4.8 and the engine study's backend settings,
including native token input. It reuses H, H+guard, PPO-1, and PPO-2
without policy search or PPO retraining. Each arm runs twice on three
constructed windows, plus one all2 anchor per window: 27 cells,
24 of which are policy cells. Lookup-69, Lookup-30, and Lookup-SG retain startup
values of 69.219, 30, and 27.568\,s and the inherited service lookup.
There are 12 unique policy/window predictions per model.

The windows are steady, middle-burst, and late-burst, with 1074, 389,
and 351 online arrivals, respectively, plus 120 offline arrivals each.
They use the existing BurstGPT-derived source and fixed thinning pipeline.
Their source rows do not overlap the canonical physical-evaluation windows
or the training/validation windows of these PPO checkpoints. Earlier
full-stream CPU replay did include these source times, and the prompt
corpus is reused: this tests additional constructed arrivals, not an
independent data source or prompt holdout.

All 27 cells pass P1+. Lookup-69 falsely rejects ten policy runs (five distinct
policy/window predictions), Lookup-30 rejects two (the two steady H repeats),
and Lookup-SG rejects none. False-accept rates are undefined because
there are no physically infeasible cells. The measured-startup model has
slightly higher total-occupation MAE than Lookup-30 (96.92 versus 94.34 GPU seconds)
but lower cost and trajectory errors overall. The trajectory improvement is
concentrated in steady; the two burst windows have slightly higher trajectory
errors than Lookup-30. Per-window and per-cell values accompany the source bundle.

No guard run in the arrival study has a risk-active tick, target raise, or reservation;
minimum queued-offline slack is 152--153\,s. These runs test fixed-parameter
prediction transfer, not an additional guard rescue.

\subsection{Startup contrast in a 14B deployment}
The 14B study uses Qwen2.5-14B-Instruct with vLLM 0.9.0, PyTorch 2.7.0+cu126,
and Transformers 4.51.1 on two A100 80GB GPUs. Each GPU hosts an independent
TP1 instance, using the main paper's FP16, eager execution, memory,
model-length, batching, and four-slot settings. The recovery and steady
input files are byte-identical to the engine study's locked inputs. H and H+guard each
run twice per window, alongside two all2 anchors per window: 12 cells,
with eight policy cells and four unique policy/window predictions per model.

Calibration uses synthetic development requests before formal prediction
or execution. The service grid crosses input lengths 128/512/2048,
output lengths 32/128/256, and concurrency 1/4, plus the 256/256,
concurrency-four offline shape. Each shape has three repeats per GPU;
the lookup takes each GPU's median batch-maximum parent-clock duration,
then the median across the two GPUs. The offline shape measures 5.844\,s.
Twenty scale-out cycles (ten per GPU) launch one instance while the other
has four admitted 256/256-token background requests. Launch-to-ready
median is 46.067\,s, range 45.814--52.352\,s. Setup is separately 48.486\,s
(maximum of two initial launches); median shutdown is 1.980\,s.

14B-Hist reuses the historical 69.219\,s startup proxy; 14B-Local uses the new
46.067\,s median. Their service, setup, shutdown, policy, and guard settings
are identical. The guard uses the measured service lookup and a separately
frozen 47.765\,s startup safety estimate (linear p95 of the 20 cycles),
with the 52.352\,s maximum for overdue-start handling and the existing 5\,s
margin. These sampling estimates carry no coverage guarantee. The source
bundle records the unrounded parameters and complete service lookup.

All 12 cells are technically valid and all four anchors pass P1+.
H leaves seven offline requests unfinished in each recovery repeat and
one in each steady repeat; all four H+guard runs satisfy P1+.
14B-Hist falsely accepts all four failed H runs. 14B-Local correctly classifies both
arms in every repeat, while lowering occupied-resource MAE from 423.35
to 68.46 and trajectory L1 from 1049.53 to 338.15 GPU seconds.
This single-field predictor contrast demonstrates the value of local startup
measurement given the shared 14B service calibration. It does not compare
reusing versus remeasuring the service curve, or isolate a causal effect of
model size, engine, or hardware.

Two guarded runs (one per window) contain 47 risk-active/reservation ticks
and one target raise each; the other two have no guard activity. The active
runs have minimum queued-offline slack of 13\,s and four/one reserved offline
dispatches in recovery/steady; inactive runs have minimum slack of 126/129\,s.
Thus the observed state trajectory, rather than the window label alone,
determines activation.

\paragraph{A fixed 30-second contrast.}
To compare local measurement with the existing simple startup assumption,
we retrospectively replay 14B-30: the same 14B service, setup, shutdown,
and guard parameters, with startup set to 30\,s. This adds four policy/window
predictions reused across the eight existing physical runs, plus two separate
all2 predictions; it adds no GPU run or parameter search. 14B-Hist and 14B-Local reproduce
their original frozen predictions. 14B-30 agrees on 6/8 policy runs, with two
false accepts (both steady H repeats) and no false rejects. Its occupied MAE,
cost MAE, and trajectory L1 are 269.12 GPU seconds, 0.2609 scenario USD, and
578.92 GPU seconds, between 14B-Hist and 14B-Local in Table~\ref{tab:e5e6}. All2 is
unchanged across the three models. This diagnoses the existing deployments;
14B-30 was not frozen before those physical runs.

\paragraph{From startup to tail deadlines.}
All three H replays first restart at 195\,s in recovery and 138\,s in steady.
At cooldown expiry (255/198\,s), 14B-Hist is still starting: its queue-plus-running
pressure is 13/14, so H raises the target to two GPUs. 14B-Local is already serving;
pressure is 5/8 and the target remains one. 14B-30 has pressure one in both
windows. These decisions alter later expansion and shutdown timing.

In recovery, physical H releases its last GPU near 1692\,s, leaving seven
later offline arrivals (1695--1785\,s) below H's scale-out threshold of nine.
They remain undispatched through the 2100\,s cutoff. In steady, the final
1785\,s offline arrival is similarly stranded after scale-down near 1780\,s.
14B-Local predicts the same failed request IDs in both repeats of each window.
14B-Hist's later shutdown lets it finish those tail requests; 14B-30 leaves six
in recovery but finishes the last steady request, explaining its false accepts.

14B-Local's guarded replay restarts capacity for the stranded queues at 1936\,s
in recovery and 2026\,s in steady. The two risk-active physical guard runs
trigger at those same ticks and complete the tail work. The other two guarded
runs finish without risk activation because their earlier trajectories differ.
Some physical-state mismatches also precede the first restart, involving
service or release timing. Thus the fixed-field replay contrast identifies
a startup feedback path; it does not attribute every physical discrepancy
to startup, or every guarded success to a guard raise. Exact event and request
chains accompany the reduced diagnostic in \path{artifact/e6_diagnostic/}.

\begin{table*}[t]
\centering
\small
\caption{Arrival study/14B policy-cell results. FA/FR denote false accepts/rejects, divided
by physically infeasible/feasible cells; a dash denotes an empty denominator.
Resource errors are GPU seconds and total-cost MAE is scenario USD, including
deadline penalties. Means weight windows equally. Repeats share predictions.}
\label{tab:e5e6}
\par\vspace{4pt}
\begin{tabular*}{\linewidth}{@{\extracolsep{\fill}}llrrrrrr@{}}
\toprule
Study & Replay & Agreement & FA & FR & Cost MAE & Occupied MAE & Trajectory L1\\
\midrule
Arrivals & Lookup-69 & 14/24 & --- & 10/24 & 0.7082 & 500.90 & 1268.65\\
Arrivals & Lookup-30 & 22/24 & --- & 2/24 & 0.0930 & 94.34 & 307.16\\
Arrivals & Lookup-SG & 24/24 & --- & 0/24 & 0.0897 & 96.92 & 250.39\\
14B & 14B-Hist & 4/8 & 4/4 & 0/4 & 0.4105 & 423.35 & 1049.53\\
14B & 14B-Local & 8/8 & 0/4 & 0/4 & 0.0668 & 68.46 & 338.15\\
\bottomrule

\end{tabular*}
\end{table*}

\paragraph{Verification and reduced evidence.}
Both campaigns finish without skipped cells or technical retries. All four
14B deadline failures remain included. CPU reanalysis reproduces the sealed
accounting and analysis files. The 14B study required a recorded identity-parser
adaptation for distribution metadata \texttt{2.7.0} versus runtime
\texttt{2.7.0+cu126}; CUDA and source identities remained checked. A prepared
external clock-recovery extension was not triggered in either campaign.
The source bundle's \path{artifact/e5_e6/} contains per-cell comparisons,
per-window summaries, guard diagnostics, and 14B calibration parameters;
\texttt{python verify.py} recomputes Table~\ref{tab:e5e6} from these reduced
comparisons. This verifies summary aggregation, not raw GPU execution.
The full machine archives remain separately retained. The arrival study/14B request and
lifecycle projections are included in the expanded paired-trace corpus;
the retrospective 14B-30 diagnostic is separate from its frozen predictions.
